% ===== PRÉAMBULE CANONIQUE — à coller VERBATIM en tête de CHAQUE .tex =====
\documentclass[11pt,a4paper]{article}
\usepackage[utf8]{inputenc}\usepackage[T1]{fontenc}\usepackage{lmodern}
\usepackage[french]{babel}
\usepackage{amsmath,amssymb,mathtools}\usepackage{icomma}
\usepackage[version=4]{mhchem}          % \ce{...} : équations & formules
\usepackage{siunitx}\sisetup{locale=FR} % \qty{}{}, \si{}, \ang{}
\usepackage{chemfig}                    % \chemfig{...} : structures développées
\usepackage{xcolor}\usepackage{colortbl}
\usepackage{tikz}\usepackage{pgfplots}\pgfplotsset{compat=1.18}
\usepackage[most]{tcolorbox}\usepackage{enumitem}\usepackage{array,booktabs}
\usepackage[a4paper,margin=2.2cm]{geometry}
\usetikzlibrary{arrows.meta,calc,angles,quotes,positioning,patterns,backgrounds,shapes.geometric}
\definecolor{primary}{RGB}{222,49,99}\definecolor{primarydark}{RGB}{150,22,55}
\definecolor{defcol}{RGB}{0,114,178}\definecolor{methcol}{RGB}{52,92,148}
\definecolor{excol}{RGB}{0,135,90}\definecolor{attncol}{RGB}{200,40,55}
\tcbset{boxbase/.style={enhanced,breakable,boxrule=0.9pt,arc=2.5pt,
  left=3mm,right=3mm,top=2.3mm,bottom=2.3mm,fonttitle=\bfseries,coltitle=white}}
% --- boîtes TUTORIEL (noms EXACTS, ne pas renommer) ---
\newtcolorbox{cle}[1][]{boxbase,colback=defcol!6,colframe=defcol,
  title={\textsc{À retenir}\ifx\relax#1\relax\relax\else\textmd{ : \itshape #1}\fi}}
\newtcolorbox{methode}[1][]{boxbase,colback=methcol!7,colframe=methcol,sharp corners,
  title={\textsc{Méthode}\ifx\relax#1\relax\relax\else\textmd{ --- \itshape #1}\fi}}
\newtcolorbox{exemplebox}[1][]{boxbase,colback=excol!5,colframe=excol,
  title={\textsc{Exemple}\ifx\relax#1\relax\relax\else\textmd{ : \itshape #1}\fi}}
\newtcolorbox{piege}[1][]{boxbase,colback=attncol!6,colframe=attncol,
  title={\textsc{Piège}\ifx\relax#1\relax\relax\else\textmd{ : \itshape #1}\fi}}
\newtcolorbox{remarque}[1][]{boxbase,colback=black!4,colframe=black!45,
  title={\textsc{Remarque}\ifx\relax#1\relax\relax\else\textmd{ : \itshape #1}\fi}}
% --- boîtes EXERCICE / CORRIGÉ (noms EXACTS, auto-comptées) ---
\newtcolorbox[auto counter]{exo}[1][]{boxbase,colback=primary!4,colframe=primary,
  title={\textsc{Exercice}~\thetcbcounter\ifx\relax#1\relax\relax\else\textmd{ --- \itshape #1}\fi}}
\newtcolorbox[auto counter]{corrige}[1][]{boxbase,colback=excol!4,colframe=excol,
  title={\textsc{Corrigé}~\thetcbcounter\ifx\relax#1\relax\relax\else\textmd{ --- \itshape #1}\fi}}
\newcommand{\R}{\mathbb{R}}\newcommand{\N}{\mathbb{N}}\newcommand{\Z}{\mathbb{Z}}\newcommand{\ds}{\displaystyle}
% (un \usepackage{fancyhdr}+en-tête est possible ; il est ignoré côté web)
% =========================================================================
\begin{document}
\sisetup{per-mode=symbol}

\begin{corrige}[mélange de trois apports d'acides forts]
\begin{enumerate}[label=\alph*)]
  \item Chaque acide fort cède intégralement son \ce{H+} :
        \[ n_1(\ce{HCl}) = 0,050\times0,10 = \qty{5,0e-3}{mol},\quad
           n_2(\ce{HBr}) = 0,050\times0,10 = \qty{5,0e-3}{mol}, \]
        \[ n_3(\ce{HCl}) = 0,050\times0,60 = \qty{3,0e-2}{mol}. \]
  \item $n(\ce{H3O+}) = 0,0050 + 0,0050 + 0,030 = \qty{4,0e-2}{mol}$ ; $V_{\text{tot}} = \qty{0,100}{L}$ :
        \[ [\ce{H3O+}] = \frac{0,040}{0,100} = \qty{0,40}{mol\per L},\qquad
           \mathrm{pH} = -\log(4{,}0\times10^{-1}) = 1 - \log 4 = 1 - 0{,}60 = \mathbf{0{,}40}. \]
\end{enumerate}
\end{corrige}

\begin{corrige}[mélange de trois bases fortes]
\begin{enumerate}[label=\alph*)]
  \item $n(\ce{NaOH}) = 0,025\times1,0 = \qty{2,5e-2}{mol}$ ;
        $n(\ce{Ba(OH)2}) = 0,025\times0,10\times\mathbf{2} = \qty{5,0e-3}{mol}$ (facteur \num{2}) ;
        $n(\ce{KOH}) = 0,050\times0,20 = \qty{1,0e-2}{mol}$.
  \item $n(\ce{OH-}) = 0,025 + 0,0050 + 0,010 = \qty{4,0e-2}{mol}$ ; $V_{\text{tot}} = \qty{0,100}{L}$ :
        \[ [\ce{OH-}] = \frac{0,040}{0,100} = \qty{0,40}{mol\per L},\quad \mathrm{pOH} = -\log(4{,}0\times10^{-1})
           = 0{,}40,\quad \mathrm{pH} = 14 - 0{,}40 = \mathbf{13{,}60}. \]
\end{enumerate}
\end{corrige}

\begin{corrige}[diluer pour tripler le degré d'ionisation]
\begin{enumerate}[label=\alph*)]
  \item $\alpha \approx \sqrt{K_a/C_a} \propto 1/\sqrt{C_a}$. Pour multiplier $\alpha$ par \num{3}, il faut diviser
        $C_a$ par $3^2 = \mathbf{9}$.
  \item La quantité de matière est conservée : diviser $C$ par \num{9} revient à multiplier le volume par \num{9}.
        On passe de \qty{1,0}{L} à $V_f = \qty{9,0}{L}$, soit \textbf{ajouter \qty{8,0}{L} d'eau}.
\end{enumerate}
\end{corrige}

\begin{corrige}[préparer une solution à partir d'un acide fort concentré]
\begin{enumerate}[label=\alph*)]
  \item La dilution conserve la quantité de matière : $n = C\,V = 5{,}0 \times 0{,}040 = \qty{0,20}{mol}$.
  \item $C_a = \dfrac{n}{V_f} = \dfrac{0,20}{1,0} = \qty{0,20}{mol\per L}$. \ce{HCl} étant fort :
        \[ \mathrm{pH} = -\log(0{,}20) = -\log(2{,}0\times10^{-1}) = 1 - \log 2 = 1 - 0{,}30 = \mathbf{0{,}70}. \]
\end{enumerate}
\end{corrige}

\begin{corrige}[capstone : même concentration, acidités différentes]
\begin{enumerate}[label=\alph*)]
  \item \ce{HCl} (fort) : $\mathrm{pH} = -\log(0{,}010) = \mathbf{2}$.
  \item \ce{HA} (faible) : $\mathrm{pH} = \tfrac12\bigl(4{,}8 - \log(10^{-2})\bigr) = \tfrac12(4{,}8 + 2)
        = \tfrac{6{,}8}{2} = \mathbf{3{,}4}$.
  \item La solution de \ce{HCl} est la plus acide (pH plus bas). Le rapport des concentrations vaut
        \[ \frac{[\ce{H3O+}]_{\ce{HCl}}}{[\ce{H3O+}]_{\ce{HA}}} = \frac{10^{-2}}{10^{-3{,}4}} = 10^{1{,}4}
           \approx 25. \]
        À concentration \emph{égale}, l'acide fort libère environ \textbf{25 fois} plus d'ions \ce{H3O+} :
        c'est toute la différence entre « fort » (totalement dissocié) et « faible ».
\end{enumerate}
\end{corrige}

\end{document}
